\documentclass[a4paper,12pt]{article}
\usepackage[utf8]{inputenc}
\usepackage[T1]{fontenc}
\usepackage[ngerman]{babel}
\usepackage{amsmath}
\usepackage{amssymb}
\usepackage{enumitem}
\usepackage{geometry}
\geometry{a4paper, margin=2.5cm}
\usepackage{parskip}

\title{Einsendeaufgaben -- Aufgabe 5.3\\[0.5em]
\large Modul 61112: Lineare Algebra}
\author{}
\date{}

\begin{document}
\maketitle

\section*{Aufgabe 5.3}

\begin{enumerate}[label=\alph*)]
  \item Es gilt für alle $n \in \mathbb{N}$:
  \begin{align*}
    (AB)^n \cdot A
      &= \underbrace{(AB \cdots AB)}_{AB:\ n\text{-mal}} \cdot A \\
      &= A \cdot \underbrace{(BA \cdots BA)}_{BA:\ n\text{-mal}} \\
      &= A \cdot (BA)^n.
  \end{align*}
  Daraus können wir schließen:
  \begin{align*}
    p(AB) \cdot A
      &= \sum_{j=0}^{n} a_j \, (AB)^j \cdot A \\
      &= \sum_{j=0}^{n} a_j \, A \cdot (BA)^j \\
      &= A \cdot \sum_{j=0}^{n} a_j \, (BA)^j \\
      &= A \cdot p(BA).
  \end{align*}

  \item Es gilt
  \[
    0 = \mu_{AB}(AB) = \mu_{AB}(AB) \cdot A = A \cdot \mu_{AB}(BA) \quad (\text{wegen a)}).
  \]
  Da $A$ invertierbar ist, folgt
  \[
    0 = A^{-1} \cdot 0 = A^{-1} \cdot A \cdot \mu_{AB}(BA) = \mu_{AB}(BA).
  \]
  Also $\mu_{AB} \in I_{BA}$ und demnach $\mathrm{grad}(\mu_{AB}) \geq \mathrm{grad}(\mu_{BA})$.
  Nach der gleichen Argumentation folgt $\mu_{BA} \in I_{AB}$ und $\mathrm{grad}(\mu_{BA}) \geq \mathrm{grad}(\mu_{AB})$.

  Wegen $\mathrm{grad}(\mu_{AB}) \geq \mathrm{grad}(\mu_{BA})$ und $\mathrm{grad}(\mu_{BA}) \geq \mathrm{grad}(\mu_{AB})$ können wir schließen
  \[
    \mathrm{grad}(\mu_{AB}) = \mathrm{grad}(\mu_{BA}).
  \]
  Da $\mu_{AB}$ das Polynom $\mu_{BA}$ erzeugt hat, existiert ein $\delta \in K[X]$ mit $\mu_{BA} = \mu_{AB} \cdot \delta$.
  Wegen $\mathrm{grad}(\mu_{AB}) = \mathrm{grad}(\mu_{BA})$ muss allerdings $\mathrm{grad}(\delta) = 0$ gelten.
  Es muss sogar $\delta = 1$ sein, da $\mu_{AB}$ und $\mu_{BA}$ normiert sind. Daraus folgt
  \[
    \mu_{BA} = \mu_{AB}.
  \]
\end{enumerate}

\end{document}
